Man rechnet jeweils beide Grössen in die SI-Einheit um.
\begin{abcliste}
\abc $\aLeO = \aLe$ und $\aReO = \aRe$, also $\aLeO \;\result{=}\; \aReO$
\abc $\bLeO = \bLe$ und $\bReO = \bRe$, also $\bLeO \;\result{=}\; \bReO$
\abc $\cLeO = \cLe$ und $\cReO = \cRe$, also $\cLeO \;\result{>}\; \cReO$
\abc $\dLeO = \dLe$, also $\dLeO \;\result{>}\; \dReO$
\abc $\eLeO = \eLe$ und $\eReO = \eRe$, also $\eLeO \;\result{=}\; \eReO$
\abc $\fLeO$ ist ein Volumen, $\fReO$ eine Masse: $\fLeO \;\result{\neq}\; \fReO$
\end{abcliste}
